\section{Check the complete cost, not just the counters}
\label{sec:cost-validation}
Correct work counts are necessary, but they do not make a latency prediction exact. Two loops
can visit the same number of rows while reading very different memory locations. A small
candidate set reused repeatedly can stay in cache; different candidates may require fresh
memory reads. The checks below separate those effects and then return to complete native calls.

\subsection{Measure buffers at the sizes the formulas describe}
The standalone profile uses the shared C++ score table and top-$K$ code. Its path operation
allocates a root bitmap, makes both children at each split, keeps the alternatives alive,
releases each parent, enumerates the final leaf and releases the remaining masks. The surviving
path and alternatives are checked to form a complete, non-overlapping partition.

Widths range from 512 bytes to 125 MB. The largest width is one bit over a billion document
positions. Thirteen stored planes and the live masks reached about 3.50 GB of process memory.
That is a measured buffer set, not the full 256-plane index. The whole-index calculation still
has to include every stored plane, packed document and ID.

The larger buffers showed median split costs of roughly 0.34--0.37 ns per physical word in this
standalone loop. At 256 dimensions, contiguous scoring took about 27 ns per document at one
and four million rows. Small postings cost more per row because preparation and selection are
less amortized. Compiler context and the surrounding query still matter; these observations
are not a universal instruction-price table.

\subsection{A scored row does not always have the same price}
A second operation reproduces the leaf-scoring loop: read a bitmap word, find its set bits,
fetch those document codes and IDs, score them, and maintain top $K$. It uses the same scorer
as the sequential loop. One run reuses a candidate bitmap after warmup. Another prepares a
different bitmap before each measured trial.

\begin{center}\small
\begin{tabular}{@{}lrr@{}}\toprule
16M-document buffer & Reused candidates & Changing candidates\\\midrule
About 1,958 scored rows & 0.159 ms & 0.455 ms\\
About 62,500 scored rows & 4.400 ms & 8.271 ms\\\bottomrule
\end{tabular}\normalsize\end{center}
The selection densities are the same; the selected IDs are not. Candidate generation is outside
the timer, but its effect on cache history remains. This shows why the scan's cost per row should
not simply be copied into B's scattered leaf scores. It does not directly measure cache-miss
counters or establish fixed CPU frequency.

A three-term fit used a fixed cost, physical bitmap words and scored rows. It was fitted on
1M and 4M pools and checked on a 16M pool. The reused-candidate model underestimated one check
by 45.2\%. With changing candidates, the worst error was still an underestimate of 27.6\%.
Both fits remain saved. The largest row pool was about 640 MB; it was not a 40 GB row-store test.

Do not add the enumeration-only leaf timer to the gathered-leaf timer: the latter already
includes enumeration. This would double-count one stage. Similarly, a fitted zero intercept
is not evidence that score-table preparation is free; other fitted terms can absorb that cost.

\subsection{Use a compact independent reference for larger native runs}
The native check uses 1M, 2M and 4M binary documents, then a separate 8M run. It keeps the
constructed strong/weak query distribution. Packed codes are generated in small chunks and
the first million rows are checked against the earlier document pool.

An independent reference avoids expanding those rows into a huge float matrix. For each query,
count the strong and weak mismatches directly from the packed bits. With magnitudes $a$ and
$\epsilon$, the score is
\[
S=a(m-2e_{\rm strong})+\epsilon(d-m-2e_{\rm weak}).
\]
Sort by this score and then ID. Small tests compare that calculation with the earlier direct
float sum, including a dimension count that does not fill the final packed word. All complete
native calls in the scaling check matched the independent reference.

For B, the leaf threshold grows with the corpus: $L=160N/10^6$. This preserves the intended
13 strong-coordinate cuts. With fixed depth and density, split words, leaf words and scored
rows all grow approximately in proportion to $N$. A simple complete-call model is then
\[
\widehat T(N)=a_0+b_0N.
\]
The fixed term covers work that does not grow with the corpus; the slope combines the work
that does. This is a conditional model of the specified path, not a reason to keep $L=128$
while increasing $N$ indefinitely.

\subsection{Two model mistakes, and how the checks exposed them}
The first fit minimized percentage error over individual queries but was used to predict a
median. Fast queries received more weight, which was the wrong target. For example, workloads
with times $[1,10,10]$ and $[2,20,20]$ have medians 10 and 20. Weighting each query by inverse
time favors the fast minority and predicts the wrong median. The repaired fit uses one median
per corpus size. A test fails for the mixed workload under the old calculation and stays correct
for a uniform workload. The original fit is preserved.

A linear row term was also insufficient for C's small postings. The shared result heap has
visible selection cost when only a few hundred documents are scored. Under independently
ordered continuous scores, document $i$ enters the current top $K$ with probability $K/i$.
The expected number of replacements after the first $K$ rows is
\[
\sum_{i=K+1}^{F}\frac{K}{i}\approx K\ln(F/K).
\]
For fixed $K$, a useful approximation is
\begin{equation}
\widehat T_C(F)=a_C+b_CF+c_C\ln\!\bigl(\max(1,F/K)\bigr).
\label{eq:small-posting-time}
\end{equation}
Binary ties and key ordering change the exact replacement count, so the logarithm is a model
term, not an exact counter. A large contiguous run supplies the row price; the native 1M/2M
measurements fit the remaining terms. This corrected the small-posting prediction without
changing the reference ranking or search implementation.

The revised 4M check is diagnostic because its results had already been inspected. The
coefficients were frozen before the 8M run. That run checks transfer to a larger corpus using
the same prescribed query workload; it is not a new-query generalization test.

\begin{center}\small
\begin{tabular}{@{}llrrr@{}}\toprule
Query support & Method & Measured ms & Predicted ms & Error\\\midrule
Fixed & A & 197.029 & 203.016 & $+3.0\%$\\
Fixed & B & 0.854 & 0.884 & $+3.6\%$\\
Fixed & C & 0.0481 & 0.0518 & $+7.7\%$\\
Changing & A & 196.664 & 197.816 & $+0.6\%$\\
Changing & B & 0.900 & 0.859 & $-4.5\%$\\
Changing & C & 202.767 & 207.086 & $+2.1\%$\\\bottomrule
\end{tabular}\normalsize\end{center}
All six checks met the earlier 20\% criterion. This supports the stated model over the checked
sizes. It does not turn the largest observed error into a guaranteed error bound at a billion
rows. The earlier failed fits remain part of the evidence.

\subsection{Check a parameter suggested by the distribution}
At one million rows, the ideal 13-bit group has mean $\lambda\approx122$ and standard deviation
about 11. A leaf size of 128 can force another cut when the strong-matching group is slightly
larger than average. A rough upper margin $\lambda+3\sqrt{\lambda}$ is about 155; choosing 160
avoids many unnecessary weak-coordinate cuts. The preceding 12-bit group has mean about 244,
so it will usually still be split. These are distribution-based settings to check, not universal
thresholds or a guarantee from a normal approximation.

The bounded follow-up adds 16,384 clusters, neighboring key widths and leaf sizes, including
160 for the changing-coordinate case. It remeasures prior choices rather than comparing fresh
timings against old ones. The first 40 generated queries become the tuning set; query IDs
40--59 are used only after the new choices are frozen. Document and reference prefixes are
checked before reusing any router.

At 1M rows and the 99\% target, the changing-coordinate follow-up selected B with one cluster
and $L=160$. It measured 0.1446 ms at 100\% binary recall. A measured 24.9235 ms at 99.95\%:
a paired speedup estimate of 172.3, with 95\% interval 169.1--177.2. C chose a 14-bit key and
measured 14.6301 ms at 100\%, but its paired speedup interval against A included one. That
point estimate alone does not establish a C advantage.

The fixed-coordinate follow-up selected a 16,384-cluster, two-probe scan and a global 13-bit
key index. Both reached 100\% recall: A measured 0.0140 ms and C 0.0128 ms. Their small gap
again concerns closely related candidate selection, not a general dominance theorem.

\subsection{Why the key method's median can change abruptly}
In the changing-support model, a fixed $h$-coordinate key captures a random number $S_h$ of the
$m$ strong coordinates. Its distribution is hypergeometric:
\begin{equation}
\Pr(S_h=s)=\frac{\binom ms\binom{d-m}{h-s}}{\binom dh}.
\label{eq:captured-strong}
\end{equation}
For $d=256$, $m=13$, the probability of capturing none is 52.75\% at $h=12$, 49.94\% at $h=13$,
and 47.27\% at $h=14$. A small width change can move a sample across the halfway point between
queries that score nearly all rows and queries that can discard about half.

Under the strong-score gap and the additional condition that weak-key bounds cannot prune
individual weak patterns, the work for a query capturing $s$ strong coordinates is approximately
\[
H=2^{h-s},\qquad F=N2^{-s}.
\]
The expected scored fraction is $\sum_s\Pr(S_h=s)2^{-s}$: about 71.30\% at $h=13$ and 69.44\%
at $h=14$. The mean work changes smoothly even when a small sample's median changes sharply.
This is another reason to keep query-level uncertainty and the full latency distribution.
The condition is about the stated constructed distribution, not arbitrary embeddings or IVF.

\subsection{A shortest-key control on new real queries}
The original real grid started with four-bit keys. A final control gives C a one-bit key and
A's routing layout, with no candidate or lookup limit. It keeps A and B's frozen settings and
uses 100 further query IDs, excluding all 500 previously used IDs. No parameter is changed
based on this control's outcomes.

A measured 9.541 ms, B 9.855 ms and C 10.545 ms. All three reached 98.95\% binary recall,
narrowly below the requested 99\%. The miss stays in the report; no tolerance is added to turn
it into a pass. There is no supported alternative speedup in this check. The original 200-query
evaluation remains a separate result, where its conservative choices met their sample targets.
A finite-sample selection rule is not a promise that every later query sample will reach the same mean.

\subsection{Measure the text query as one request}
The final pipeline check runs actual text through the cached local Nomic model on MPS, applies
the same normalization, routes the resulting vector and calls the native search. Batch size is
one. Each sequence shape is warmed first; model loading, index building and reference creation
are outside the request timer. References are computed afterward from the actual encoded vectors.

\begin{center}\small
\begin{tabular}{@{}lrrrr@{}}\toprule
Method & Encode median & Retrieve median & Total median & Total p95\\\midrule
A & 6.488 ms & 9.406 ms & 16.313 ms & 19.201 ms\\
B & 6.500 ms & 9.651 ms & 16.303 ms & 19.217 ms\\
C & 7.543 ms & 11.201 ms & 18.912 ms & 21.898 ms\\\bottomrule
\end{tabular}\normalsize\end{center}
These are 20 existing independent-query IDs, three blocks and two repeats, with the earlier
frozen settings. All reached 99.45\% binary recall on this sample. The tiny A/B difference does
not establish a speedup. Each total is measured per request: its median is not the sum of the
two stage medians. The process lifetime peak was about 3.67 GB including reference construction;
MPS driver allocation was about 1.08 GB and may overlap RSS, so those numbers are not simply added.

This text check applies to the real Nomic queries. The constructed strong-coordinate vectors
do not come from this encoder, so its time must not be added to those examples and called a
measured end-to-end sparse-embedding pipeline.
